\uuid{OC7P}
\titre{Ensembles}
\auteur{bodin ; croizat ; sacre}
\organisation{exo7}
\licence{CC BY-NC-SA 4.0}
\datecreate{2022-03-08}
\niveau{L1}
\module{Algèbre}
\chapitre{Logique, ensemble, raisonnement}
\sousChapitre{Ensemble}
\difficulte{2}
\theme{Ensembles}
\exo7qcmid{hilisit.logique.030}

\contenu{
\begin{question}
\qclassification{100.02}
Soient $f,g,h$ des fonctions définies par
$$f(x) = x^2+x \qquad g(x) = 2x-1 \qquad h(x) = \frac{1}{x+1}.$$
Quelles sont les affirmations vraies ?
\begin{answers}
	\good{$(f \circ g)(x) =4x^2-2x$}
	\bad{$(g \circ f)(x) = 2x^2+2x$}
	\good{$(h \circ f)(x) = \frac 1{x^2+x+1}$}
	\bad{$(g \circ h)(x) = \frac{1}{2x}$}
\end{answers}
\begin{explanations}
$(f \circ g)(x) = f(g(x)) = (2x-1)^2+(2x-1) = 4x^2-2x$. \\
$(g \circ f)(x) = g(f(x)) = 2(x^2+x)-1 = 2x^2+2x-1$. \\
$(h \circ f)(x) = h(f(x)) = \frac{1}{f(x)+1}=\frac 1{x^2+x+1}$.\\
$(g \circ h)(x) = g(h(x)) = 2\frac{1}{x+1} -1 = \frac{2}{x+1} - \frac{x+1}{x+1} = \frac{1-x}{x+1}$. Ce qui est proposé est $(h\circ g)(x)$.
\end{explanations}
\end{question}
}
